\begin{table}[htbp]
\centering
\caption{Model-implied dispensary-revenue benchmark, Colorado border counties, 2024}
\label{tab:counterfactual}
\footnotesize
\begin{tabular}{lc}
\toprule
Quantity & Value \\
\midrule
Estimated interaction coefficient $\theta$ & $-0.458$ \\
Free-entry benchmark multiplier $e^{-\theta}$ & 1.58 \\
\midrule
Actual 2024 revenue, CO border counties & \$324.2 M \\
Free-entry-benchmark counterfactual revenue (NM) & \$512.5 M \\
Model-implied additional dispensary revenue & +\$188.3 M \\
95\% county-cluster bootstrap interval & (\$-71.0 M, +\$899.5 M) \\
\bottomrule
\end{tabular}
\begin{minipage}{0.98\textwidth}
\vspace{0.5em}

\fontsize{9pt}{11pt}\selectfont
\textit{Notes:} The benchmark multiplier $e^{-\theta}$ is applied to each
Colorado border county's own 2024 dispensary revenue and summed, rather than to
one pre-aggregated total; New Mexico is the free-entry benchmark state and
contributes only to the estimation of $\theta$. Colorado is restricted to 2022Q2
onward (Section~\ref{sec:crossstate}). The bootstrap interval resamples counties
with replacement within each of the two states (500 replicates), refitting the
full interaction regression each time; it is wide and comfortably includes zero,
because $\theta$ is not statistically distinguishable from zero once Colorado's
pre-2022 years are dropped. The dollar figure should therefore be read as an
illustrative benchmark rather than an identified estimate of foregone revenue.
\end{minipage}
\end{table}
